\chapter{Касательная геометрия ранга один и идеал вырожденных двух форм}

Предыдущая глава вернула пятимерному пространству
\(\rho M_3+M_3\rho\) правильный статус операторного угла, но условие
\(QXQ=0\) оставалось внешним. Проверим одновременно его геометрическое
происхождение и судьбу кривизны после факторизации по идеалу вырожденных форм.

\section{Условие опоры как касательное пространство}

Пусть
\begin{equation}
 \rho=\operatorname{diag}(1,0,0),\qquad Q=I-\rho.
\end{equation}
Многообразие матриц ранга не выше единицы задаётся всеми минорами второго
порядка. Линеаризация этих девяти уравнений в точке \(\rho\) имеет ранг
четыре, поэтому её ядро имеет комплексную размерность пять. Прямой расчёт
даёт
\begin{equation}
 T_\rho\{X:\rank X\leq1\}
 =\{Z:QZQ=0\}
 =\rho M_3+M_3\rho.
 \label{eq:v5-rank-one-tangent}
\end{equation}
Тем самым прежнее условие опоры не требуется вставлять отдельно: оно есть
касательное условие аффинного конуса ранга один.

Для фиксированной амплитуды радиальное направление \(\mathbb C\rho\)
отделяется, а касательное пространство орбиты равно
\begin{equation}
 \mathcal E_\rho=\rho M_3Q\oplus QM_3\rho,
 \qquad \dim_{\mathbb C}\mathcal E_\rho=4.
 \label{eq:v5-rank-one-orbit-tangent}
\end{equation}

\section{Минимальная алгебра нулевой степени}

Стабилизатор разложения \(\mathbb C^3=\rho\mathbb C^3\oplus Q\mathbb C^3\)
задаёт
\begin{equation}
 \mathcal A_\rho=\mathbb C\rho\oplus QM_3(\mathbb C)Q
 \simeq\mathbb C\oplus M_2(\mathbb C).
 \label{eq:v5-rank-one-zero-algebra}
\end{equation}
Берём произвольный ненулевой соединитель
\begin{equation}
 D_w=\begin{pmatrix}0&w^*\\w&0\end{pmatrix},
 \qquad w\in\mathbb C^2\setminus\{0\}.
\end{equation}
Представленные одноформы \(a[D_w,b]\) заполняют в точности
\(\mathcal E_\rho\). Для любого \(w\ne0\) существуют \(r>0\) и
\(U\in U(2)\), для которых \(w=rUe_1\). Оператор
\(\operatorname{diag}(1,U)\) нормализует \(\mathcal A_\rho\), а ненулевое
масштабирование \(D_w\) не меняет пространств форм. Поэтому достаточно
координатного соединителя, для которого ранг одноформ равен четырём;
остальные образцы служат проверкой устойчивости. Радиальная амплитуда не
потеряна: она относится к нулевой степени, тогда как одноформы описывают
переходы между опорным и дополнительным углами.

\section{Факторизация двух форм}

Вторая степень до факторизации имеет вид
\begin{equation}
 \pi(\Omega_u^2\mathcal A_\rho)
 =\mathbb C\rho\oplus QM_3Q,
 \qquad \dim_{\mathbb C}=5.
\end{equation}
Для каждого соотношения
\begin{equation}
 \sum_{a,b}c_{ab}\,a[D_w,b]=0
\end{equation}
его дифференциал даёт вырожденную двухформу
\begin{equation}
 \sum_{a,b}c_{ab}[D_w,a][D_w,b].
\end{equation}
Унитарная эквивалентность и машинный аудит на координатном и четырёх общих
соединителях устанавливают точные равенства пространств
\begin{equation}
 \pi(d\ker\pi_1)=QM_3Q,
 \qquad
 \Omega_{D_w}^2(\mathcal A_\rho)
 \simeq\mathbb C\rho.
 \label{eq:v5-rank-one-twoform-quotient}
\end{equation}
Следовательно, новый минимальный носитель не повторяет полностью прежний
восемнадцатимерный запрет: на физическом углу остаётся одна самосопряжённая
двухформа. Однако вся матричная форма дополнительного двумерного сектора
удаляется идеалом вырожденных форм.

\section{Проверка ориентации}

Пусть \(d_w\) содержит только переход из \(\rho\)-угла в \(Q\)-угол. Тогда
\begin{align}
 [d_w,d_w^*]&=\operatorname{diag}(-\|w\|^2,ww^*),\\
 D_w^2&=\operatorname{diag}(\|w\|^2,ww^*).
\end{align}
После факторизации общий блок \(ww^*\) исчезает и остаётся
\begin{equation}
 [d_w,d_w^*]\longmapsto-\|w\|^2\rho,
 \qquad
 D_w^2\longmapsto+\|w\|^2\rho.
\end{equation}
Их знаки различны, но квадраты норм совпадают. Значит, обычное положительное
действие на факторе снова не различает ориентации и не восстанавливает
матричное отображение момента.

\begin{theorem}[Касательный успех и запрет кривизны формы]
Ранг-один проектор канонически порождает условие \(QXQ=0\), а его
стабилизаторная алгебра воспроизводит внедиагональный связывающий модуль как
пространство одноформ. Но обычный фактор двух форм по идеалу вырожденных
форм одномерен: он сохраняет только \(\mathbb C\rho\) и уничтожает
\(QM_3Q\). Поэтому эта геометрия выводит радиальную опорную кривизну, но не
полное ориентированное отображение момента.
\end{theorem}

\begin{protocol}
\begin{enumerate}
 \item происхождение условия \(QXQ=0\) из ранга один: \textbf{получено};
 \item связывающий угол как пространство одноформ: \textbf{получен};
 \item устойчивость рангов относительно выбора ненулевого \(w\):
 \textbf{проверена};
 \item фактор двух форм: \textbf{одномерная опорная линия};
 \item сохранение матричной кривизны \(QM_3Q\): \textbf{опровергнуто};
 \item выбор ориентации квадратом нормы: \textbf{не получен};
 \item физическое замыкание: \textbf{не достигнуто}.
\end{enumerate}
\end{protocol}

Следующая проверка должна выяснить, может ли уже построенный коммутирующий
квадрат наблюдаемого сектора превратить исчезнувший \(Q\)-блок в
\emph{относительную} бимодульную кривизну. Повторять обычный фактор Конна на
том же носителе больше не следует.

Машинный сертификат сохранён в
\path{s2t_v5_rank_one_tangent_junk_gate.py} и
\path{s2t_v5_rank_one_tangent_junk_gate_results.json}.